% Section 1: shared introduction for the report.
% The method chapters (Sections 2, 3 and 4) stay unchanged.

\section{Introduction}

This project studies face reenactment, one way to make a deepfake. The aim is to
take the face from one picture and let it follow the movement from another
source, then write the result as a video. The assignment asked for a working
deepfake model, or any model with a comparable face reenactment or
identity-mapping effect. It also asked us to document every choice and turning
point, to include the failed experiments with the reasons for the failure, and
to reflect on what each experiment taught us. This report gives that record.

Every method in this report works with two inputs. The identity image supplies
the face. The driving or motion source supplies the movement. The two roles
never change, but each method uses the pair its own way, and the chapters give
the inputs their own names. The autoencoder chapter calls them the source face
image and the target video. The FOMM chapter calls them the source image and the
driving video. The rectified flow chapter pairs a portrait with a body frame
from the clip.

We explored three methods, one from each family. The first is an autoencoder
with one shared encoder and one decoder per person. We built and trained it
ourselves to learn the basics, and we also tested three pretrained face-swap
models from FaceFusion. The second is the keypoint-based First Order Motion
Model (FOMM). It needs no landmark detector and no 3D face model, and it
animates a source photo with the motion of a driving video. The third is a head
swap with a rectified flow model and a LoRA adapter. A released checkpoint that
needs only four sampling steps fits the one laptop with the 4 GB GPU that we
had.

The three methods differ inside, but they share one picture-level pipeline. They
read an identity image and a motion source, run a model that writes one output
frame per frame of the motion source, and save the frames as a clip. The
autoencoder had to be trained first. FOMM and the rectified flow pipeline load
released checkpoints, and the rectified flow pipeline adds restoration,
smoothing and reassembly after generation. We judge the frames by eye and with
the measures defined in each chapter.

Section 2 covers the autoencoder experiments and the pretrained FaceFusion
models, Section 3 covers the GAN background and FOMM, and Section 4 covers the
rectified flow pipeline. Section 5 compares the three methods side by side, and
the conclusion chapter reports what worked, what failed, and what we learned.
